\section*{Ringraziamenti}
\paragraph{Camilla Maccari} mi sono appoggiato molto sui
suoi esercizi svolti nei vari tentativi di superare l'esame. I
seguenti sono poco pi\'u di una rifinitura e qualche completamento
delle note di Camilla.

\paragraph{Vittorio Scirli} lo ringrazio molto per avermi fatto scopire
il software \textit{falstad}, che ha permesso di simulare i circuiti con
facilit\'a e di dare un criterio di compatibilit\'a tra i risultati
da me ottenuti e la ``realt\'a'' (simulata da falstad). Lo ringrazio
anche per avermi aiutato in generale con questa materia.

\paragraph{Marco Nencioni} lo ringrazio per essere stato il mio
compagno di laboratorio dal \textit{primo} fino al \textit{terzo}
anno. Siamo passati, da non saper fare il fit di una retta, a non capire
la differenza fra serie e parallelo fino a non capire come fare un
semaforo.

\section*{Leggimi}
I seguenti esercizi non sono guidati dal lato di vista teorico, il
loro scopo \'e quello di dare un riferimento durante la preparazione
dell'esame alla luce della teoria \textit{gi\'a imparata} da materiale
esterno a questo; tuttavia pu\'o essere anche utile questo testo come
fonte di esempi per accompagnare l'apprendimento.
Detto questo buona fortuna!
